\documentclass[11pt,a4paper]{article}

\usepackage[utf8]{inputenc}
\usepackage[russian,english]{babel}
\usepackage[T2A]{fontenc}
\usepackage{amsmath,amssymb,amsthm}
\usepackage{hyperref}
\usepackage{geometry}
\geometry{margin=2.5cm}

\title{%
Alanian-GRA Multiverse:\\
Многоуровневая GRA Мета-обнулёнка\\
и аланский когнитивный вакуум
}

\author{AAA AAA}
\date{2026}

\begin{document}

\maketitle

\begin{abstract}
В работе предлагается архитектура Alanian-GRA Multiverse --- специализированный вариант
многоуровневой GRA Мета-обнулёнки, калиброванный на аланско-осетинский когнитивный код.
Аланская традиция (честь, гостеприимство, клановая вложенность, трифункциональная структура
Нартского эпоса) рассматривается как этно-онтологический тест-стенд для мультиверсных
когнитивных архитектур. Формализуются пространства состояний, цели и проекторы,
определяется функционал мультиверса и категория ``аланских когнитивных вакуумов'',
в которой аланский вакуум выступает терминальным объектом. Это обеспечивает
культурно-насыщенное понятие объектности и согласованности уровней, потенциально
пригодное для проектирования сложных AI-систем.
\end{abstract}

\section{Введение}

Многоуровневая архитектура GRA Мета-обнулёнки описывает иерархию состояний и целей,
согласуемых через рекурсивное обнуление ``пены'' несогласованности между доменами
и уровнями. Такая архитектура стремится к \emph{когнитивному вакууму} --- состоянию,
в котором артефакты интерпретации сняты на всех уровнях иерархии, а сохраняются только
структурные инварианты.

В данной работе предлагается специальный кейс этой архитектуры --- Alanian-GRA
Multiverse, в котором \emph{инварианты} задаются аланско-осетинским когнитивным кодом:
честью, гостеприимством, клановой вложенностью и трифункциональной структурой,
реконструированной по Нартскому эпосу.

\section{Многоуровневая GRA Мета-обнулёнка}

Кратко приведём формальную схему базовой GRA-мультиверс-архитектуры. Пусть
\(\mathbf{a} = (a_0, a_1, \dots, a_k)\) --- мультииндекс, задающий положение
подсистемы в иерархии уровней, и \(\mathcal{H}^{(\mathbf{a})}\) --- соответствующее
пространство состояний. Полное пространство мультиверса:
\[
\mathcal{H}_{\text{multiverse}} =
\bigotimes_{l=0}^{K} \mathcal{H}^{(l)}, \quad
\mathcal{H}^{(l)} = \bigotimes_{\dim(\mathbf{a})=l} \mathcal{H}^{(\mathbf{a})}.
\]

Для каждого уровня \(l\) задана цель \(G_l\) и проектор \(\mathcal{P}_{G_l}\).
Пена уровня \(l\) определяется как
\[
\Phi^{(l)}(\Psi^{(l)}, G_l) =
\sum_{\mathbf{a}\neq \mathbf{b} \atop \dim(\mathbf{a})=\dim(\mathbf{b})=l}
\big|\langle \Psi^{(\mathbf{a})} \mid \mathcal{P}_{G_l} \mid \Psi^{(\mathbf{b})} \rangle\big|^2.
\]

Полный функционал мультиверса:
\[
J_{\text{multiverse}}(\mathbf{\Psi}) =
\sum_{l=0}^{K} \Lambda_l
\sum_{\dim(\mathbf{a})=l} J^{(l)}(\Psi^{(\mathbf{a})}),
\]
где \(J^{(l)}\) определяется рекурсивно, а веса \(\Lambda_l = \lambda_0 \alpha^l\),
\(0<\alpha<1\), задают затухание по уровням.

\section{Аланская когнитивная архитектура}

Аланско-осетинская традиция рассматривается как конкретная реализация
многоуровневой когнитивной архитектуры. На уровне индивида и ситуации
фиксируются роли воина, хозяина, гостя и родича, а также три функции
(суверенно-сакральная, воинская, хозяйственно-социальная), согласующиеся
с трифункциональной схемой Ж.~Дюмезиля.

Клановые, племенные и союзные уровни формируют иерархию, в которой решения
индивидов тензорно собираются в конфигурации родов, общин, союзов и царства
Алании. Высшие уровни включают имперско-религиозный и общекавказский контекст.

\section{Инвариантные проекторы Alanian-GRA}

В Alanian-GRA вводятся три жёстких инварианта в виде проекторов:

\begin{itemize}
  \item \(\mathcal{P}_{\text{honor}}\): отбрасывает компоненты состояний, нарушающие
  фундаментальные клятвы, честное слово и базовый код чести.
  \item \(\mathcal{P}_{\text{hospitality}}\): при активном гостевом контексте
  переводит гостя в подпространство полной защиты, запрещая агрессию и
  монетизацию помощи.
  \item \(\mathcal{P}_{\text{kin}}\): сохраняет структурную целостность клановой
  и родовой вложенности.
\end{itemize}

Эти проекторы предполагается коммутирующими и действующими на всех уровнях
иерархии, определяя аланские цели \(G_l^{\text{Alan}}\) и их проекторы
\(\mathcal{P}_{G_l^{\text{Alan}}}\).

\section{Функционал Alanian-мультиверса и обнуление}

Alanian-мультиверсный функционал
\[
J_{\text{Alan-multiverse}}(\mathbf{\Psi}) =
\sum_{l=0}^{K} \Lambda_l
\sum_{\dim(\mathbf{a})=l} J^{(l)}(\Psi^{(\mathbf{a})})
\]
строится так, чтобы пены \(\Phi^{(l)}\) минимизировались при одновременном
соблюдении инвариантов чести, гостеприимства и клановости.

Формулируется теорема мультиверсного обнуления для аланского случая:
при коммутативности проекторов и согласованности иерархии существует
состояние \(\Psi^*_{\text{Alan}}\), для которого \(\Phi^{(l)} = 0\) для всех
уровней, а инварианты выполняются автоматически.

\section{Категория аланских когнитивных вакуумов}

На языке теории категорий определяем подкатегорию AlanCogVac, объекты которой
— аланские архитектуры \((\mathcal{H}_{\text{Alan}}, \mathcal{P}, \mathcal{I}_{\text{Alan}},
\mathcal{F}_{\text{tri}})\), а морфизмы — отображения, сохраняющие иерархию
проекторов и инвариантные проекторы.

Аланский когнитивный вакуум \(\Psi^*_{\text{Alan}}\) задаёт терминальный
объект категории: для любой аланской архитектуры существует единственный
морфизм, стягивающий её в это состояние, интерпретируемый как операция
полного обнуления пены при сохранении инвариантов.

\section{Заключение}

Alanian-GRA Multiverse соединяет строгий формализм многоуровневой GRA
Мета-обнулёнки с конкретной этно-онтологической традицией. Это позволяет
переосмыслить объектность и согласованность уровней как \emph{верность
культурным инвариантам}, а не абстрактную нейтральность, и даёт шаблон
для проектирования мультиверсных AI-систем с жёстко заданным
когнитивным ядром.

\end{document}
